\section{Планируемые результаты работы}

В результате выполнения работы будет разработан и выложен воспроизводимый экспериментальный пайплайн для VLM-guided генерации и редактирования изображений диффузионными моделями. Пайплайн будет включать реализацию оптимизации не только начального латента, но и промежуточных латентных состояний в процессе денойзинга, а также непрерывных управляющих параметров генерации (в частности, коэффициента classifier-free guidance и параметров, используемых в инверсии при редактировании, включая null-text эмбеддинги).

Будет реализован единый протокол экспериментов для text-to-image и image editing режимов, включающий подготовку наборов промптов различной сложности и стандартизированную процедуру генерации с фиксированным сидом, числом шагов денойзинга и бюджетом вычислений для сравнения методов.

Будет проведено количественное сравнение предложенного подхода с базовыми генераторами (SD1.5, SD3.0, FLUX.1[dev], FLUX.2[dev]) и с proxy-prompts подходом SAP \cite{huberman2025contextually}. Сравнение будет выполнено по метрикам семантического соответствия: CLIP Score (по протоколу из статьи Noise Diffusion \cite{miao2024noisediffusion}) и alignment на основе VLM (по протоколу из статьи SAP \cite{huberman2025contextually}) на наборах промптов разной сложности.

Будут получены отдельные результаты, характеризующие вклад каждого компонента предложенного подхода: оптимизации только $z_T$, оптимизации набора промежуточных $z_t$, совместной оптимизации $z_t$ и коэффициента CFG, а также (для редактирования) совместной оптимизации $z_t$ и null-text эмбеддингов. Для сравнения будет использован единый вычислительный бюджет, что позволит корректно сопоставить эффективность различных вариантов.

Будет подготовлен набор качественных результатов (визуальные примеры) для промптов с нетипичными композиционными и количественными ограничениями.

По итогам экспериментов будут сформулированы выводы о применимости VLM-gradient guidance для выхода за пределы prior диффузионных моделей в задачах генерации и редактирования, а также практические рекомендации по настройке пайплайна (какие латенты оптимизировать, на каких шагах, с каким шагом оптимизации и какие параметры управления дают наибольший эффект при заданном бюджете).